\documentclass[11pt]{article}

\usepackage[a4paper,margin=1in]{geometry}
\usepackage[T1]{fontenc}
\usepackage[utf8]{inputenc}
\usepackage{lmodern}
\usepackage{amsmath,amssymb,mathtools}
\usepackage{booktabs}
\usepackage{float}
\usepackage{xcolor}
\usepackage{listings}
\usepackage{hyperref}

\hypersetup{
  colorlinks=true,
  linkcolor=blue!55!black,
  citecolor=blue!55!black,
  urlcolor=blue!60!black
}

\definecolor{codebg}{RGB}{247,248,250}
\definecolor{codecomment}{RGB}{42,112,74}
\definecolor{codekeyword}{RGB}{42,74,145}
\definecolor{codestring}{RGB}{150,66,55}

\lstdefinelanguage{Wolfram}{
  morekeywords={ClearAll,Needs,DiracTrace,DiracSimplify,DiracSigma,
    Contract,ScalarProductExpand,ExpandScalarProduct,Collect2,Simplify,
    Together,Thread,Table,Association,AssociationThread,True,False},
  sensitive=true,
  morecomment=[s]{(*}{*)},
  morestring=[b]"
}

\lstset{
  language=Wolfram,
  basicstyle=\ttfamily\small,
  keywordstyle=\color{codekeyword},
  commentstyle=\color{codecomment},
  stringstyle=\color{codestring},
  backgroundcolor=\color{codebg},
  frame=single,
  rulecolor=\color{black!15},
  breaklines=true,
  columns=fullflexible,
  keepspaces=true,
  showstringspaces=false,
  aboveskip=0.8em,
  belowskip=0.8em
}

\newcommand{\ii}{\mathrm{i}}
\newcommand{\ee}{\mathrm{e}}
\newcommand{\Tr}{\operatorname{Tr}}
\newcommand{\slashedk}{\not\!k}
\newcommand{\slashedp}{\not\!p}
\newcommand{\slashedpp}{\not\!p^{\prime}}
\newcommand{\slashedq}{\not\!q}
\newcommand{\Mathematica}{\textit{Mathematica}}
\newcommand{\FeynCalc}{\textsc{FeynCalc}}
\newcommand{\GeV}{\mathrm{GeV}}
\newcommand{\keV}{\mathrm{keV}}

\title{Independent Mathematica Derivation for\\
  $D_{s1}/B_{s1}\to D_s/B_s\,\gamma$}
\author{Ds/Bs gamma project}
\date{24 July 2026}

\begin{document}
\maketitle

\begin{abstract}
This note is a clean, independent restart of the symbolic part of the
$D_{s1}/B_{s1}\to D_s/B_s\gamma$ light-cone sum-rule calculation.  It records
exactly what has been implemented and checked in \Mathematica{} and
\FeynCalc: the currents, correlator ordering, Wick contraction, propagator
routing, Rohrwild three-point/background-field notation, hard-emission
traces, Ward identities, E1 projection,
inverse-denominator reduction, Feynman parameterization, the first soft
two-particle Fierz traces, physical-state mixing, the pure-axial Colangelo
limit, the decay-constant normalization, the complete axial and tensor
post-double-Borel invariants, and the numerical couplings and widths.  The
Mathematica numerical implementation is independent of the Python
implementation; their output tables are compared only after both
calculations have finished.  The 24--26 July revision separates off-shell
QCD kinematics from on-shell pole kinematics and withdraws the earlier
physical-mass substitution in the theoretical tensor numerator.
\end{abstract}

\tableofcontents

\section{Purpose and executable files}

The interactive source of this calculation is
\begin{center}
  \texttt{notebooks/DsBs\_gamma\_symbolic\_derivation.nb}.
\end{center}
It contains 53 input cells designed to be evaluated from top to bottom or one
at a time.  The noninteractive regression scripts are
\begin{align*}
  &\texttt{scripts/mathematica\_current\_wick\_trace\_audit.wl},\\
  &\texttt{scripts/mathematica\_correlator\_routing\_audit.wl},\\
  &\texttt{scripts/mathematica\_hard\_loop\_parameterization.wl},\\
  &\texttt{scripts/mathematica\_state\_mixing\_projection.wl},\\
  &\texttt{scripts/step14\_explicit\_double\_borel\_forms.wl},\\
  &\texttt{scripts/step15\_complete\_three\_particle\_borel.wl}.
\end{align*}
The first three batch audits write their full machine-readable expressions
under \texttt{outputs/}; the state-mixing audit prints its exact residuals and
interval overlaps directly.  This note presents the compact analytic content
needed to understand and reproduce those expressions.

\section{Conventions and kinematics}

We use
\begin{equation}
  g_{\mu\nu}=\operatorname{diag}(1,-1,-1,-1),\qquad
  \{\gamma_\mu,\gamma_\nu\}=2g_{\mu\nu},
\end{equation}
\begin{equation}
  \sigma_{\mu\nu}=\frac{\ii}{2}[\gamma_\mu,\gamma_\nu],\qquad
  \epsilon^{0123}=+1.
\end{equation}
The \FeynCalc{} convention is explicitly checked through
\texttt{FeynCalc`\$LeviCivitaSign=-1}.

The momentum assignment is
\begin{equation}
  p^{\prime}=p+q,
\end{equation}
where $p^{\prime}$ is the initial $1^+$ momentum, $p$ is the final $0^-$
momentum, and
$q$ is the photon momentum.  The photon is real,
\begin{equation}
  q^2=0,\qquad 2p\cdot q=p^{\prime 2}-p^2.
\end{equation}
The heavy flavor is denoted by $Q=c,b$ and its mass by $m_Q$; the light flavor
is strange with mass $m_s$.

\section{Currents and the Rohrwild correlator}

The two basis-current vertices are
\begin{align}
  \Gamma^A_\mu
  &=\gamma_\mu\gamma_5,\\
  \Gamma^B_\mu
  &=\frac{\ii}{m_Q+m_s}\,
    \sigma_{\mu\alpha}p^{\prime\alpha}\gamma_5,
\end{align}
and the pseudoscalar vertex is
\begin{equation}
  \Gamma_P=\ii\gamma_5.
\end{equation}
The tensor vertex is independently represented as
\begin{equation}
  \Gamma^B_\mu
  =-\frac{(\gamma_\mu\slashedpp-\slashedpp\gamma_\mu)\gamma_5}
    {2(m_Q+m_s)}.
\end{equation}
\Mathematica{} verifies that the two definitions agree exactly.

The physical-current rotations are kept outside the basis traces:
\begin{align}
  \Gamma^{\rm low}_\mu
  &=\sin\theta\,\Gamma^A_\mu+\cos\theta\,\Gamma^B_\mu,\\
  \Gamma^{\rm high}_\mu
  &=\cos\theta\,\Gamma^A_\mu-\sin\theta\,\Gamma^B_\mu.
\end{align}

Following Rohrwild's vacuum three-point notation
\href{https://arxiv.org/abs/0708.1405}{arXiv:0708.1405}, adapted to a
$1^+\to0^-\gamma$ transition, the starting correlator is
\begin{equation}
  \Pi^X_{\mu\nu}(p,q)=\ii^2\int d^4x\,d^4y\,
  \ee^{\ii p\cdot x+\ii q\cdot y}
  \langle0|T\{J_P^\dagger(x)j_\nu^{\rm em}(y)J^X_\mu(0)\}|0\rangle,
  \qquad X=A,B,
  \label{eq:rohrwild-correlator}
\end{equation}
where
\begin{equation}
  j_\nu^{\rm em}=e_s\bar s\gamma_\nu s+e_Q\bar Q\gamma_\nu Q.
\end{equation}
Here $e_s,e_Q$ are charges in units of the positive elementary charge $e$.
Introduce the plane-wave background field
\begin{equation}
  F_{\alpha\beta}(z)=\ii
  (q_\alpha\varepsilon^*_\beta-q_\beta\varepsilon^*_\alpha)
  \ee^{\ii q\cdot z}.
\end{equation}
The term linear in this field obeys
\begin{equation}
  \varepsilon^{*\nu}\Pi^X_{\mu\nu}(p,q)
  =\ii\int d^4x\,\ee^{\ii p\cdot x}
  \langle0|T\{J_P^\dagger(x)J^X_\mu(0)\}|0\rangle_F
  \equiv
  \ii\int d^4x\,\ee^{\ii p\cdot x}
  \langle\gamma(q,\varepsilon)|T\{J_P^\dagger(x)J^X_\mu(0)\}|0\rangle,
  \label{eq:background-external-equivalence}
\end{equation}
up to the explicitly stated overall QED convention.  Thus the former
vacuum-to-photon formula was not a different calculation; it is the
single-photon component of Eq.~\eqref{eq:rohrwild-correlator}.

In explicit fields the ordered product is
\begin{equation}
  T\{\bar s(x)\Gamma_PQ(x)\,\bar Q(0)\Gamma^X_\mu s(0)\}.
  \label{eq:field-order}
\end{equation}
Thus the vertex at $x$ carries the final channel $p$, while the current at the
origin carries the initial channel $p^{\prime}=p+q$.

The current definitions used in the notebook are:
\begin{lstlisting}
gammaA[index_] := GA[index] . GA[5];
gammaB[index_] :=
  I FV[pPrime, alpha] DiracSigma[GA[index], GA[alpha]] . GA[5]/
    (mQ + ms);
gammaBExpanded[index_] :=
  -(GA[index] . GS[pPrime] - GS[pPrime] . GA[index]) . GA[5]/
    (2 (mQ + ms));
gammaP[] := I GA[5];
propNumerator[momentum_, mass_] := GS[momentum] + mass;
\end{lstlisting}

\section{Wick contraction and global signs}

For the complete hard contraction, Eq.~\eqref{eq:field-order} gives
\begin{equation}
  \langle T\{\bar s\Gamma_PQ\,\bar Q\Gamma_Xs\}\rangle
  =-\Tr_D[\Gamma_PS_Q\Gamma_XS_s].
\end{equation}
The minus sign is the closed-fermion-loop sign.  Including the leading
correlator factor $\ii$, but not yet the propagator or QED-vertex factors, the
hard prefactor is therefore
\begin{equation}
  -\ii N_c.
\end{equation}

If only the heavy flavor is contracted, the strange bilocal remains open and
there is no closed fermion loop.  The soft contraction begins with $+\ii$.
The Fierz rearrangement
\begin{equation}
  \bar s_\alpha(x)s_\beta(0)
  =-\frac14\sum_i(\Gamma_i)_{\beta\alpha}
  \bar s(x)\Gamma_i s(0)
\end{equation}
then gives a soft prefactor $-\ii/4$ before the free-propagator factor.

In the notebook the bookkeeping is kept visible:
\begin{lstlisting}
fermionLoopSign = -1;
hardWickPrefactor = I fermionLoopSign Nc;
softWickPrefactorBeforeFierz = I;
softWickPrefactorAfterFierz = -I/4;
\end{lstlisting}

\section{Propagator expansion and OPE building blocks}

The free momentum-space propagator is
\begin{equation}
  S_f^0(r)=\frac{\ii(\not\!r+m_f)}{r^2-m_f^2+\ii0}.
  \label{eq:free-propagator}
\end{equation}

\subsection{Light-strange propagator in Rohrwild notation, including $m_s$}

To make the comparison with Rohrwild direct, write $\bar u=1-u$ and keep the
background fields at the point $ux$.  With color indices displayed, the
light-strange propagator used as the master expansion is
\begin{align}
  \widehat{s^a(x)\bar s^b(0)}={}&
  \delta^{ab}\Bigg[
    \frac{\ii\not\!x}{2\pi^2x^4}
    -\frac{m_s}{4\pi^2x^2}
  \Bigg]
  \nonumber\\
  &-\frac{\ii g_s}{16\pi^2x^2}\int_0^1du\,
  \left\{\bar u\not\!x\sigma_{\alpha\beta}
  +u\sigma_{\alpha\beta}\not\!x\right\}
  G_A^{\alpha\beta}(ux)t^A_{ab}
  \nonumber\\
  &-\frac{\ii e_s\delta^{ab}}{16\pi^2x^2}\int_0^1du\,
  \left\{\bar u\not\!x\sigma_{\alpha\beta}
  +u\sigma_{\alpha\beta}\not\!x\right\}
  F^{\alpha\beta}(ux)
  \nonumber\\
  &-\frac{m_s}{32\pi^2}\int_0^1du\,
  \left[g_sG_A^{\alpha\beta}(ux)t^A_{ab}
       +e_s\delta^{ab}F^{\alpha\beta}(ux)\right]
  \sigma_{\alpha\beta}L_x+\cdots ,
  \label{eq:light-full-propagator}
\end{align}
where
\begin{equation}
  L_x=\ln\!\left(\frac{-x^2\mu^2}{4}\right)+2\gamma_E .
\end{equation}
The first line contains only the free propagator and its explicit $m_s$ term.
The last line is the linear-$m_s$ Wilson coefficient in a single gluon or
electromagnetic background field.  No local vacuum-condensate term is written
in Eq.~\eqref{eq:light-full-propagator}.  Setting $m_s$ and the last line to
zero gives precisely the light-propagator form displayed in Eq.~(2.12) of
Rohrwild.  The ellipsis denotes higher-background-field and higher-order mass
terms, not local condensate pieces used elsewhere.

\paragraph{Important calculation warning.}
Although local parts are no longer displayed as part of the propagator, the
existing axial numerical sum-rule implementation \emph{does include} the
term labeled \texttt{heavy\_local} in the Python audit.  After the double
Borel transform it is
\begin{equation}
  \Pi^{A,\mathrm{local}}_{Q\gamma}
  =e_Q\,\ee^{-m_Q^2/M^2}\langle\bar ss\rangle
  \left[
    1-\frac{m_Qm_s}{M^2}
    +\frac{m_s^2}{2M^2}\left(1+\frac{m_Q^2}{M^2}\right)
  \right].
  \label{eq:heavy-local}
\end{equation}
It has not been removed here, because doing so would change the quoted form
factors and widths.  The mixed condensate
$\langle\bar s g_s\sigma\cdot Gs\rangle$ and the vacuum gluon condensate were
not included in the radiative three-point calculation.

\subsection{Heavy propagator in a background gluon field}

For reference, the momentum representation through one background gluon can
be organized as
\begin{align}
  S_Q^{ab}(x)
  =\ii\int\frac{d^4k}{(2\pi)^4}\,\ee^{-\ii k\cdot x}
  \Bigg\{&
  \frac{\delta^{ab}(\not\!k+m_Q)}{k^2-m_Q^2}
  -\frac{g_sG^A_{\alpha\beta}t^A_{ab}}{4}
  \frac{
    \sigma^{\alpha\beta}(\not\!k+m_Q)
    +(\not\!k+m_Q)\sigma^{\alpha\beta}}
  {(k^2-m_Q^2)^2}
  +\cdots\Bigg\}.
  \label{eq:heavy-full-propagator}
\end{align}
For the photon-DA calculation the one-gluon term is used as an operator
insertion, rather than replaced by $\langle G^2\rangle$.  The explicit
light-cone form used in the symbolic three-particle kernels is
\begin{align}
  S_Q^G(x)
  =-\ii g_s\int\frac{d^4k}{(2\pi)^4}\,\ee^{-\ii k\cdot x}
  \int_0^1du\Bigg[
  &\frac{\not\!k+m_Q}{2(m_Q^2-k^2)^2}
    G^{\alpha\beta}(ux)\sigma_{\alpha\beta}
  \nonumber\\
  &+\frac{ux_\alpha}{m_Q^2-k^2}
    G^{\alpha\beta}(ux)\gamma_\beta
  \Bigg].
  \label{eq:gluon-background-propagator}
\end{align}
Both the $\sigma_{\alpha\beta}$ part and the
$x_\alpha\gamma_\beta$ part have Mathematica trace-kernel scripts.

\subsection{Electromagnetic background-field propagator}

The first-order electromagnetic correction is
\begin{align}
  S_f^\gamma(x)
  =-\ii e_f\int\frac{d^4k}{(2\pi)^4}\,\ee^{-\ii k\cdot x}
  \int_0^1du\Bigg[
  &\frac{\not\!k+m_f}{2(m_f^2-k^2)^2}
    F^{\alpha\beta}(ux)\sigma_{\alpha\beta}
  \nonumber\\
  &+\frac{ux_\alpha}{m_f^2-k^2}
    F^{\alpha\beta}(ux)\gamma_\beta
  \Bigg].
  \label{eq:em-background-propagator}
\end{align}
The overall electromagnetic coupling may either be included in $e_f$ or kept
outside the propagator, but not both.  In the explicit-vertex trace audit,
Eq.~\eqref{eq:em-background-propagator} is represented equivalently by
\begin{equation}
  S_f^0(r)(-\ii e e_f\gamma_\nu\varepsilon^\nu)S_f^0(r+q),
\end{equation}
with the orientation fixed separately for each diagram.  These are alternative
representations of the same hard photon-emission term.

\subsection{Photon-DA matrix elements}

The condensate-dependent photon terms are vacuum-to-photon nonlocal matrix
elements.  We use Rohrwild's background-field notation directly:
\begin{align}
  \langle0|\bar s(0)[0,x]\sigma_{\mu\nu}s(x)|0\rangle_F
  ={}&e e_s\langle\bar ss\rangle\int_0^1du\,
  \Bigg\{
  \left[\chi_\gamma\phi_\gamma(u)
        +\frac{x^2}{16}\mathcal A(u)\right]F_{\mu\nu}(ux)
  \nonumber\\
  &\hspace{0.2cm}
  +\frac{\mathcal B(u)}{8}x^\rho
  \left[x_\nu F_{\mu\rho}(ux)-x_\mu F_{\nu\rho}(ux)\right]
  \Bigg\},
  \label{eq:tensor-photon-da}
\end{align}
and
\begin{equation}
  \langle0|\bar s(0)[0,x]\gamma_\mu s(x)|0\rangle_F
  =-\frac{e e_s}{2}f_{3\gamma}
  \int_0^1du\,\overline\psi^{(V)}(u)x^\rho F_{\rho\mu}(ux).
  \label{eq:vector-photon-da}
\end{equation}
Thus no $H_\gamma$ or $\bar H_\gamma$ is introduced in the Rohrwild
formulae.  The exact notation dictionary to the alternative $h_\gamma$
representation is
\begin{equation}
 \boxed{\mathcal B(u)
 =-4\int_0^u d\alpha\,(u-\alpha)h_\gamma(\alpha)},\qquad
 \mathbb A(u)-8\bar H_\gamma(u)
 =\mathcal A(u)+2\mathcal B(u).
 \label{eq:rohrwild-B-dictionary}
\end{equation}
The factor two in the second equality is essential.  It follows from the
operator identity in the Ball--Braun--Kivel convention; replacing the
Colangelo combination by $\mathcal A+\mathcal B$ would be incorrect.
Similarly,
\begin{equation}
 \overline\psi^{(V)}(u)=2\int_0^u dv\,\psi^v(v)
 =2\Psi^v(u)
 \label{eq:rohrwild-vector-dictionary}
\end{equation}
for the conformal models used here.  The three-particle matrix elements
\[
 \langle\gamma|\bar s(x)g_sG_{\alpha\beta}(vx)\Gamma s(0)|0\rangle
\]
define $\mathcal S,\widetilde{\mathcal S},\mathcal T_1,\ldots,\mathcal T_4$.
Their axial combination used in the sum rule is
\begin{equation}
  \mathcal F_1=
  \mathcal S+\widetilde{\mathcal S}-\mathcal T_1-\mathcal T_2
  +\mathcal T_3+\mathcal T_4
  +2v(-\mathcal S-\mathcal T_3+\mathcal T_2).
  \label{eq:F1-combination}
\end{equation}

\section{Momentum routing}

\subsection{Heavy-line emission}

The three fermion momenta are $k$, $k+q$, and $k-p$.  The denominators are
\begin{align}
  d_{H1}&=(k+q)^2-m_Q^2,\\
  d_{H2}&=k^2-m_Q^2,\\
  d_{H3}&=(k-p)^2-m_s^2.
\end{align}
The current momenta are
\begin{equation}
  k-(k-p)=p,\qquad (k+q)-(k-p)=p+q.
\end{equation}

\subsection{Strange-line emission}

The three momenta are $k$, $k-p$, and $k-p-q$.  The denominators are
\begin{align}
  d_{S1}&=k^2-m_Q^2,\\
  d_{S2}&=(k-p)^2-m_s^2,\\
  d_{S3}&=(k-p-q)^2-m_s^2.
\end{align}
The corresponding current momenta are
\begin{equation}
  k-(k-p)=p,\qquad k-(k-p-q)=p+q.
\end{equation}

This vertex check is important.  The old trial route $k-p+q$ gives
$k-(k-p+q)=p-q$ at the initial vertex and is therefore not the routing of
Eq.~\eqref{eq:rohrwild-correlator}.  It is kept only in the dedicated routing
audit as a diagnostic comparison.

\section{Hard-emission Dirac traces}

Charges, color, denominators, the QED vertex, propagator factors, and the
global Wick sign are deliberately kept outside the numerator traces.  For a
generic basis vertex $\Gamma^X_\mu$, the heavy-emission numerator is
\begin{equation}
  N^X_{H,\mu\nu}=\Tr_D\!\left[
  \ii\gamma_5(\slashedk+m_Q)\gamma_\nu
  (\slashedk+\slashedq+m_Q)\Gamma^X_\mu
  (\slashedk-\slashedp+m_s)\right].
  \label{eq:heavy-trace}
\end{equation}
The strange-emission numerator is
\begin{equation}
  N^X_{s,\mu\nu}=\Tr_D\!\left[
  \ii\gamma_5(\slashedk+m_Q)\Gamma^X_\mu
  (\slashedk-\slashedp-\slashedq+m_s)\gamma_\nu
  (\slashedk-\slashedp+m_s)\right].
  \label{eq:strange-trace}
\end{equation}

The implementation is short enough to display in full:
\begin{lstlisting}
traceReduce[expression_] :=
  Collect2[
    DiracSimplify[expression, DiracSubstitute67 -> True],
    {mQ, ms}
  ];

heavyTrace[current_] := traceReduce[DiracTrace[
  gammaP[] . propNumerator[k, mQ] . GA[nu] .
  propNumerator[k + q, mQ] . current[mu] .
  propNumerator[k - p, ms]
]];

strangeTrace[current_] := traceReduce[DiracTrace[
  gammaP[] . propNumerator[k, mQ] . current[mu] .
  propNumerator[k - p - q, ms] . GA[nu] .
  propNumerator[k - p, ms]
]];

nAHeavy   = heavyTrace[gammaA];
nAStrange = strangeTrace[gammaA];
nBHeavy   = heavyTrace[gammaB];
nBStrange = strangeTrace[gammaB];
\end{lstlisting}
Each of the four traces is recomputed after a cyclic rotation of the matrix
chain.  All four cyclic differences vanish exactly.

\section{Fermion-line Ward identities}

For equal masses on the photon-emitting segment and $B=A+q$,
\begin{equation}
  (\not\!A+m)\not\!q(\not\!B+m)
  =D_B(\not\!A+m)-D_A(\not\!B+m),
  \label{eq:ward-line}
\end{equation}
where $D_A=A^2-m^2$ and $D_B=B^2-m^2$.  Equation~\eqref{eq:ward-line}
shows explicitly that replacing the photon vertex by $q_\nu$ collapses one of
the adjacent propagators, as required by the Ward--Takahashi identity.

The notebook checks Eq.~\eqref{eq:ward-line} independently for the heavy line,
with $A=k$ and $B=k+q$, and the strange line, with $A=l-q$ and $B=l$:
\begin{lstlisting}
wardHeavyResidual = FCE[DiracSimplify[
  DiracGammaExpand[
    propNumerator[k, mQ] . GS[q] . propNumerator[k + q, mQ]
    - (SP[k + q, k + q] - mQ^2) propNumerator[k, mQ]
    + (SP[k, k] - mQ^2) propNumerator[k + q, mQ]
  ],
  DiracSubstitute67 -> True, DiracOrder -> True
]] // ScalarProductExpand // Simplify;
\end{lstlisting}
Both the heavy and strange residuals are zero.

\section{E1 projection}

The transverse E1 tensor is chosen as
\begin{equation}
  S_{\mu\nu}=p_\nu q_\mu-(p\cdot q)g_{\mu\nu}.
\end{equation}
For $q^2=0$,
\begin{equation}
  q^\nu S_{\mu\nu}=0.
\end{equation}
The normalized projector is
\begin{equation}
  \mathcal P_{E1}^{\mu\nu}
  =\frac{p^\nu q^\mu-(p\cdot q)g^{\mu\nu}}
  {2(p\cdot q)^2},
\end{equation}
and satisfies
\begin{equation}
  \mathcal P_{E1}^{\mu\nu}S_{\mu\nu}=1.
\end{equation}
The code is
\begin{lstlisting}
e1Tensor = FV[p, nu] FV[q, mu] - SP[p, q] MT[mu, nu];
e1Projector = e1Tensor/(2 SP[p, q]^2);

e1Project[expression_] := Collect2[
  ScalarProductExpand[
    Contract[expression e1Projector] /. kinematicRules
  ],
  {mQ, ms}
];

projAHeavy   = e1Project[nAHeavy];
projAStrange = e1Project[nAStrange];
projBHeavy   = e1Project[nBHeavy];
projBStrange = e1Project[nBStrange];
\end{lstlisting}
The projector normalization, photon Ward contraction, and contraction with a
physical initial-state polarization are all checked symbolically.

\section{Physical-state projection and the Colangelo limit}

\subsection{Rotation of the projected invariant amplitudes}

Let $\widehat T_A$ and $\widehat T_B$ denote the continuum-subtracted,
double-Borel E1 invariant amplitudes obtained from the two basis currents.
Because the correlator is linear in the current at the initial-state vertex,
the physical projections are
\begin{equation}
  \begin{pmatrix}\widehat T_1\\[2pt]\widehat T_2\end{pmatrix}
  =R(\theta)
  \begin{pmatrix}\widehat T_A\\[2pt]\widehat T_B\end{pmatrix},
  \qquad
  R(\theta)=
  \begin{pmatrix}
    \sin\theta&\cos\theta\\
    \cos\theta&-\sin\theta
  \end{pmatrix},
  \qquad R^TR=\mathbf 1.
  \label{eq:state-rotation}
\end{equation}
The same rotation therefore applies after the E1 projection, after the OPE,
and after the Borel transformation; these linear operations commute.

Define the pseudoscalar and physical axial-state residues by
\begin{align}
  \langle0|J_P|P(p)\rangle
  &=\frac{f_Pm_P^2}{m_Q+m_s},
  \label{eq:pseudoscalar-residue}\\
  \langle0|J_i^\mu|H_i(p^{\prime},\eta)\rangle
  &=f_iM_i\eta^\mu,
  \qquad i=1,2,
  \label{eq:physical-decay-constant-definition}
\end{align}
and use
\begin{equation}
  \langle\gamma(q,\varepsilon)P(p)|H_i(p^{\prime},\eta)\rangle
  =eG_i\big[(\varepsilon^*\!\cdot\eta)(p\cdot q)
  -(\varepsilon^*\!\cdot p)(\eta\cdot q)\big].
  \label{eq:mixed-transition-amplitude}
\end{equation}
The pole term in the projected, double-Borel invariant is then
\begin{equation}
  \widehat T_i^{\rm pole}
  =\frac{f_Pm_P^2}{m_Q+m_s}\,f_iM_iG_i
  \exp\!\left(-\frac{M_i^2}{\mathcal M_1^2}
               -\frac{m_P^2}{\mathcal M_2^2}\right).
  \label{eq:mixed-hadronic-pole}
\end{equation}
Consequently the correctly normalized physical coupling is
\begin{equation}
  G_i=
  \frac{m_Q+m_s}{f_Pm_P^2f_iM_i}
  \exp\!\left(\frac{M_i^2}{\mathcal M_1^2}
              +\frac{m_P^2}{\mathcal M_2^2}\right)
  \widehat T_i^{\rm OPE}.
  \label{eq:mixed-coupling-from-ope}
\end{equation}
If the basis residues for a given pole are defined with the same mass
normalization, Eq.~\eqref{eq:mixed-coupling-from-ope} becomes
\begin{align}
  G_1
  &=\frac{\sin\theta\,f_A^{(1)}G_A^{(1)}
          +\cos\theta\,f_B^{(1)}G_B^{(1)}}{f_1},
  \label{eq:residue-weighted-low}\\
  G_2
  &=\frac{\cos\theta\,f_A^{(2)}G_A^{(2)}
          -\sin\theta\,f_B^{(2)}G_B^{(2)}}{f_2}.
  \label{eq:residue-weighted-high}
\end{align}
Thus a bare sine--cosine rotation of already normalized couplings is valid
only when the corresponding residues cancel.

\subsection{Pole-side versus QCD-side kinematics}

The distinction between the two sides of the sum rule is essential.  Before
the Borel transformation the QCD correlator depends on two independent
virtualities:
\begin{equation}
 q^2=0,\qquad p^{\prime}=p+q,\qquad
 2p\cdot q=p^{\prime2}-p^2,\qquad
 p^2\ \hbox{and}\ p^{\prime2}\ \hbox{independent}.
 \label{eq:qcd-offshell-kinematics}
\end{equation}
One must \emph{not} replace these variables by hadron masses inside the
theoretical numerator.  The mass substitutions occur only in the residue of
an isolated hadronic pole,
\begin{equation}
 \left.p^2\right|_{\rm pole}=m_P^2,\qquad
 \left.p^{\prime2}\right|_{\rm pole}=M_i^2,\qquad
 \left.p\cdot q\right|_{\rm pole}=\frac{M_i^2-m_P^2}{2}.
 \label{eq:hadronic-pole-kinematics}
\end{equation}
Likewise, $u_0$ is not a pre-Borel kinematic replacement.  It is generated by
the double Borel transform and becomes $1/2$ only after the equal-Borel choice.
This corrects the former ``physical-residue numerator prescription'' used in
the provisional Python tensor estimate.

\subsection{Rohrwild decomposition before the Borel transformation}

In the Rohrwild organization the two basis invariants have parallel
decompositions,
\begin{align}
 T_A^{\rm RW}(p^{\prime2},p^2)
 &=
 T_A^{\rm pert}+T_A^{\rm tw2}+T_A^{\rm tw3}+T_A^{\rm tw4}
 +T_A^{\rm 3p,g}+T_A^{\rm 3p,\gamma},
 \label{eq:TA-preborel-decomposition}\\
 T_B^{\rm RW}(p^{\prime2},p^2)
 &=
 T_B^{\rm pert}+T_B^{\rm tw2}+T_B^{\rm tw3}+T_B^{\rm tw4}
 +T_B^{\rm 3p,g}+T_B^{\rm 3p,\gamma}.
 \label{eq:TB-preborel-decomposition}
\end{align}
There is no separate ordinary local-condensate term in either equation.
Factors of $\langle\bar ss\rangle$ still occur in the normalization of
\emph{nonlocal} photon matrix elements such as
$\chi_\gamma\langle\bar ss\rangle\phi_\gamma$; those factors must not be
confused with insertion of the local propagator term
$-\langle\bar ss\rangle/12$.

The perturbative contribution has the off-shell double-dispersion form
\begin{equation}
 T_X^{\rm pert}(p^{\prime2},p^2)
 =\int ds_1\,ds_2\,
 \frac{\rho_X^{\rm pert}(s_1,s_2)}
 {(s_1-p^{\prime2})(s_2-p^2)},\qquad X=A,B.
 \label{eq:preborel-double-dispersion}
\end{equation}
For a two-particle photon DA define
\begin{equation}
 D(u)=m_Q^2-\bar u\,p^{\prime2}-u\,p^2
     =m_Q^2-(p+\bar u q)^2,\qquad \bar u=1-u.
 \label{eq:two-particle-denominator}
\end{equation}
Then every two-particle term can be kept explicitly off shell as
\begin{equation}
 T_X^{\rm 2p}
 =\sum_{\Phi,n}\int_0^1du\,
 \frac{C_{X,\Phi}^{(n)}(u,p^2,p\cdot q)\,\Phi(u)}
 {[D(u)]^n},
 \quad
 \Phi\in\{\phi_\gamma,\overline\psi^{(V)},\mathcal A,\mathcal B\}.
 \label{eq:preborel-two-particle-master}
\end{equation}
For a three-particle DA the analogous denominator is
$D(\xi)=m_Q^2-\xi p^{\prime2}-(1-\xi)p^2$, where $\xi$ is the photon momentum
fraction fixed by $(\alpha_i,v)$, and
\begin{equation}
 T_X^{\rm 3p}
 =\sum_{\mathcal F,n}\int\mathcal D\alpha_i\int_0^1dv\,
 \frac{C_{X,\mathcal F}^{(n)}
 (\alpha_i,v,p^2,p\cdot q)\,\mathcal F(\alpha_i)}
 {[D(\xi)]^n}.
 \label{eq:preborel-three-particle-master}
\end{equation}

The current-dependent E1 trace kernels already obtained independently in
\Mathematica{} are listed below.  Common charges, Wilson denominators, and DA
normalizations are not included in the table.  For the soft two-particle
line, $k=p+\bar u q$ has been used, but no hadronic mass substitution has
been made.
\begin{table}[H]
 \centering
 \small
 \begin{tabular}{p{0.25\textwidth}p{0.30\textwidth}p{0.34\textwidth}}
  \toprule
  DA structure & $A$-current kernel & $B$-current kernel\\
  \midrule
  two-particle tensor
   & $8$
   & $\displaystyle\frac{8m_Q}{m_Q+m_s}$\\[4pt]
  two-particle vector
   & $\displaystyle-\frac{6\ii m_Q}{p\cdot q}$
   & $\displaystyle-\frac{2\ii}{m_Q+m_s}
      \frac{3p^2+(2\bar u+3)p\cdot q}{p\cdot q}$\\[7pt]
  $\mathcal S_\gamma$, $\sigma_{\alpha\beta}$ part
   & $8$
   & $\displaystyle\frac{8m_Q}{m_Q+m_s}$\\[4pt]
  $\mathcal T_4^\gamma$, $\sigma_{\alpha\beta}$ part
   & $-16\ii$
   & $\displaystyle-\frac{16\ii m_Q}{m_Q+m_s}$\\[4pt]
  $\mathcal S_\gamma$, $x_\alpha\gamma_\beta$ part
   & $\displaystyle-\frac{4\ii\,q\cdot x}{p\cdot q}$
   & $0$\\[7pt]
  $\mathcal T_4^\gamma$, $x_\alpha\gamma_\beta$ part
   & $0$
   & $0$\\
  \bottomrule
 \end{tabular}
 \caption{Explicit pre-Borel E1 kernels.  These are QCD-side expressions and
 contain the independent invariants $p^2$ and $p\cdot q$ rather than hadron
 masses.  The zeroes in the last column follow only after using the
 numerator-free $x_\alpha\mathcal F^{\alpha\beta}\gamma_\beta$ propagator
 term.}
 \label{tab:explicit-preborel-kernels}
\end{table}

\subsection{Exact double-Borel identities}

Define
\begin{equation}
 \frac{1}{M^2}=\frac{1}{\mathcal M_1^2}
                +\frac{1}{\mathcal M_2^2},\qquad
 u_0=\frac{\mathcal M_1^2}{\mathcal M_1^2+\mathcal M_2^2},
 \qquad
 A_n=\frac{(M^2)^{2-n}}{\Gamma(n)}
     \exp\!\left(-\frac{m_Q^2}{M^2}\right).
 \label{eq:double-borel-definitions}
\end{equation}
For
$I_n[F]=\int_0^1du\,F(u)/[D(u)]^n$, the double Borel transform is
\begin{equation}
 \mathcal B_{\mathcal M_1^2}^{p^{\prime2}}
 \mathcal B_{\mathcal M_2^2}^{p^2} I_n[F]
 =A_nF(u_0).
 \label{eq:basic-double-borel}
\end{equation}
For $n>1$ and vanishing endpoint terms, numerator virtualities give
\begin{align}
 \mathcal B_{12}\!\int_0^1du\,
 \frac{(p\cdot q)F(u)}{D(u)^n}
 &=\frac{A_{n-1}}{2(n-1)}F'(u_0),
 \label{eq:borel-pq}\\
 \mathcal B_{12}\!\int_0^1du\,
 \frac{p^{\prime2}F(u)}{D(u)^n}
 &=m_Q^2A_nF(u_0)-A_{n-1}F(u_0)
   +\frac{A_{n-1}}{n-1}[uF(u)]'_{u=u_0},
 \label{eq:borel-pprime2}\\
 \mathcal B_{12}\!\int_0^1du\,
 \frac{p^2F(u)}{D(u)^n}
 &=m_Q^2A_nF(u_0)-A_{n-1}F(u_0)
   -\frac{A_{n-1}}{n-1}[\bar uF(u)]'_{u=u_0}.
 \label{eq:borel-p2}
\end{align}
Here $\mathcal B_{12}$ abbreviates the two Borel operators in
Eq.~\eqref{eq:basic-double-borel}.  If an endpoint does not vanish, its contact
term must be retained.  Equations~\eqref{eq:borel-pq}--\eqref{eq:borel-p2}
show explicitly why replacing $p^2$ and $p\cdot q$ by pole masses on the QCD
side is not valid: the correct Borel images contain DA derivatives and
possible endpoint terms.

The perturbative term transforms instead according to
\begin{equation}
 \widehat T_X^{\rm pert}
 =\int_{\Omega(s_0)}ds_1\,ds_2\,
 \rho_X^{\rm pert}(s_1,s_2)
 \exp\!\left(-\frac{s_1}{\mathcal M_1^2}
             -\frac{s_2}{\mathcal M_2^2}\right).
 \label{eq:borel-double-dispersion}
\end{equation}
Thus a numerator $p^{\prime2}$, $p^2$, or $p\cdot q$ in a genuine double
dispersion density maps to $s_1$, $s_2$, or $(s_1-s_2)/2$, respectively,
up to single-pole contact terms annihilated by the second Borel transform.

\subsection{Explicit invariants after the Borel transformation}

For the equal-Borel choice
$\mathcal M_1^2=\mathcal M_2^2=2M^2$, one has $u_0=1/2$.  In the
Rohrwild momentum routing used below it is convenient to write
\begin{equation}
 a_0=1-u_0=\frac{\mathcal M_2^2}
 {\mathcal M_1^2+\mathcal M_2^2},\qquad a_0=\frac12
 \quad\hbox{for equal Borel masses}.
 \label{eq:a0-definition}
\end{equation}
The physical normalization factor and couplings are
\begin{align}
 \mathcal N_i(M^2)
 &=\frac{\exp[(M_i^2+m_P^2)/(2M^2)](m_Q+m_s)}
 {M_if_i\,m_P^2f_P},\qquad i=1,2,
 \label{eq:physical-normalization-factor}\\
 G_1&=\mathcal N_1
   \left(\sin\theta\,\widehat T_A^{\rm RW}
        +\cos\theta\,\widehat T_B^{\rm RW}\right),
 \label{eq:paper-final-g1}\\
 G_2&=\mathcal N_2
   \left(\cos\theta\,\widehat T_A^{\rm RW}
        -\sin\theta\,\widehat T_B^{\rm RW}\right).
 \label{eq:paper-final-g2}
\end{align}
The pole masses in $\mathcal N_i$ come from the phenomenological side and do
not constitute a QCD-side substitution.

In Rohrwild notation the Borel invariants are
\begin{align}
 \widehat T_A^{\rm RW}
 &=\widehat T_A^{\rm pert}+\widehat T_A^{\rm tw2}
  +\widehat T_A^{\rm tw3}+\widehat T_A^{\rm tw4}
  +\widehat T_A^{\rm 3p,g}+\widehat T_A^{\rm 3p,\gamma},
 \label{eq:basis-ope-decomposition}\\
 \widehat T_B^{\rm RW}
 &=\widehat T_B^{\rm pert}+\widehat T_B^{\rm tw2}
  +\widehat T_B^{\rm tw3}+\widehat T_B^{\rm tw4}
  +\widehat T_B^{\rm 3p,g}+\widehat T_B^{\rm 3p,\gamma}.
 \label{eq:tensor-ope-decomposition}
\end{align}
There is no $\widehat T_X^{\langle\bar ss\rangle}$ slot.

The completed axial pieces not affected by the disputed on-shell shortcut are
\begin{align}
 \widehat T_A^{\rm pert}
 &=\int_{(m_Q+m_s)^2}^{s_0}ds\,
   \ee^{-s/M^2}\rho_A^P(s),
 \nonumber\\
 \widehat T_A^{\rm tw2}
 &=e_s\langle\bar ss\rangle(E_Q-E_0)M^2
   \chi_\gamma\phi_\gamma(a_0),
 \nonumber\\
 \widehat T_A^{\rm tw3}
 &=e_sf_{3\gamma}m_QE_Q\overline\psi^{(V)}(a_0),
 \nonumber\\
 \widehat T_A^{\rm tw4}
 &=\frac{e_s\langle\bar ss\rangle E_Q}{4}
   [\mathcal A(a_0)+2\mathcal B(a_0)]
   \left(1+\frac{m_Q^2}{M^2}\right),
 \nonumber\\
 \widehat T_A^{\rm 3p,g}
 &=e_s\langle\bar ss\rangle E_Q\,
   \mathcal I_{\mathcal F_1}(u_0),
 \label{eq:full-axial-borel-invariant}
\end{align}
where
\begin{equation}
 E_Q=\ee^{-m_Q^2/M^2},\qquad E_0=\ee^{-s_0/M^2},
 \label{eq:borel-exponent-shorthand}
\end{equation}
The continuum prescription is term-specific.  The leading twist-2 term
carries $\Delta E_Q=E_Q-E_0$; twist-3, twist-4, and the gluonic
three-particle contribution retain $E_Q$, as in the channel-specific
$D_{s1}\to D_s\gamma$ sum rule.  The separate gauge-completion
electromagnetic three-particle term below follows Rohrwild's $I_F$-type
$\Delta E_Q$ subtraction.  The three-particle convolution is
\begin{align}
 \mathcal I_{\mathcal F_1}(u_0)={}&
 \int_0^{1-u_0}dv
 \int_0^{u_0/(1-v)}d\alpha_g\,
 \mathcal F_1(u_0-(1-v)\alpha_g,
              1-u_0-v\alpha_g,\alpha_g)
 \nonumber\\
 &+\int_{1-u_0}^{1}dv
 \int_0^{(1-u_0)/v}d\alpha_g\,
 \mathcal F_1(u_0-(1-v)\alpha_g,
              1-u_0-v\alpha_g,\alpha_g).
 \label{eq:F1-integral-compact}
\end{align}
The condensate factors in these equations belong to nonlocal photon DAs.

The fully assembled gluonic and electromagnetic three-particle terms are
given below in Eqs.~\eqref{eq:complete-gluonic-three-particle} and
\eqref{eq:complete-em-three-particle}.  In particular, no unevaluated
$\mathcal B_{12}[\,\cdots\,]$ remains in the final expressions.  The
previously printed terms proportional to $m_P^2-M_i^2$ applied
Eq.~\eqref{eq:hadronic-pole-kinematics} inside a QCD numerator and are
withdrawn.

For the tensor current, the perturbative density used in the diagonal
double-dispersion test is
\begin{align}
 \rho_B^P(s)&=e_s\mathcal L(s;m_s,m_Q;a_s,b_s)
              -e_Q\mathcal L(s;m_Q,m_s;a_Q,b_Q),\nonumber\\
 \mathcal L(s;m_i,m_j;a,b)&=-\frac{3}{8\pi^2}\left[
 2a\ln\frac{s-m_j^2+m_i^2-\sqrt{\lambda}}
 {s-m_j^2+m_i^2+\sqrt{\lambda}}
 +\frac{b}{m_i^2}\frac{m_j^2-m_i^2-s}{s^2}\sqrt{\lambda}\right],
 \label{eq:tensor-perturbative-density}
\end{align}
\begin{align}
 d&=m_Q+m_s,\qquad
 a_s=\frac{m_Qm_s}{d},\qquad
 b_s=\frac{m_s^2s}{d},\nonumber\\
 a_Q&=\frac{2m_Q^2-m_Qm_s}{d},\qquad
 b_Q=-\frac{m_Q^2s}{d}.
\end{align}
The tensor-DA traces imply the exact, kinematics-independent relations
\begin{equation}
 \widehat T_B^{\rm tw2}=\frac{m_Q}{d}\widehat T_A^{\rm tw2},
 \qquad
 \widehat T_B^{\rm tw4}=\frac{m_Q}{d}\widehat T_A^{\rm tw4}.
 \label{eq:tensor-safe-two-particle-ratios}
\end{equation}
The vector DA contains $p^2$ and $p\cdot q$ before the transform.  With
\begin{equation}
 D_a=m_Q^2-a p^{\prime2}-(1-a)p^2,
 \label{eq:Da-definition}
\end{equation}
its independently reduced off-shell kernels are
\begin{align}
 K_A^{(V)}(a)&=-\frac{8m_Q}{D_a^2},&
 K_B^{(V)}(a)&=-\frac{8[m_Q^2+(1-a)p\cdot q]}{dD_a^2}.
 \label{eq:vector-preborel-reduced-kernels}
\end{align}
The exact $n=2$ double-Borel maps in this variable are
\begin{align}
 \mathcal B_{12}\int_0^1da\,\frac{F(a)}{D_a^2}
 &=E_QF(a_0),\nonumber\\
 \mathcal B_{12}\int_0^1da\,\frac{(p\cdot q)F(a)}{D_a^2}
 &=-\frac{M^2E_Q}{2}F'(a_0),\nonumber\\
 \mathcal B_{12}\int_0^1da\,\frac{p^2F(a)}{D_a^2}
 &=E_Q\left[m_Q^2F(a_0)+M^2a_0F'(a_0)\right].
 \label{eq:explicit-n2-borel-maps-a}
\end{align}
Consequently the requested explicit twist-3 tensor result is
\begin{equation}
 \boxed{
 \widehat T_B^{\rm tw3}
 =\frac{e_sf_{3\gamma}E_Q}{d}
 \left[
 m_Q^2\overline\psi^{(V)}(a_0)
 -\frac{M^2}{2}
 \left.\frac{d}{da}
 \{(1-a)\overline\psi^{(V)}(a)\}\right|_{a=a_0}
 \right].}
 \label{eq:tensor-vector-explicit-borel}
\end{equation}
For the symmetric choice, $\overline\psi^{(V)}(1/2)=0$, and hence
\begin{equation}
 \widehat T_A^{\rm tw3}=0,\qquad
 \widehat T_B^{\rm tw3}
 =-\frac{e_sf_{3\gamma}M^2E_Q}{4d}
 \overline\psi^{(V)\prime}\!\left(\frac12\right).
 \label{eq:symmetric-vector-borel}
\end{equation}
Thus the vanishing of the axial vector-DA term at $a_0=1/2$ must not be
copied to the tensor current.

\paragraph{Corrected post-Borel
$x_\alpha\mathcal F^{\alpha\beta}\gamma_\beta$ kernels.}
For a two- or three-particle DA projected onto the heavy-line momentum
fraction $a$, define its effective line density
\begin{equation}
 F_{\rm eff}(a)=
 \int\mathcal D\alpha_i\int_0^1dv\,
 w_F(\alpha_i,v)\mathcal F(\alpha_i)
 \delta[a-a(\alpha_i,v)].
 \label{eq:effective-line-density}
\end{equation}
The weight $w_F$ contains the explicit $v$ factor and the Lorentz weight of
the DA.  This propagator term has no $(\slashedk+m_Q)$ numerator.
Writing $F_0=F_{\rm eff}(a_0)$ and
$F_1=F_{\rm eff}'(a_0)$, the complete scalar Borel images of the
Mathematica $x_\alpha\gamma_\beta$ E1 kernels are
\begin{align}
 \mathfrak X_{A,S}[F]&=-8E_QF_0,\qquad
 \mathfrak X_{A,T_1}[F]=48\ii E_QF_0,\nonumber\\
 \mathfrak X_{A,T_2}[F]&=16\ii E_QF_0,\qquad
 \mathfrak X_{A,T_3}[F]=-16\ii E_QF_0,\nonumber\\
 \mathfrak X_{A,T_4}[F]&=0,
 \label{eq:axial-xgamma-postborel}\\
 \mathfrak X_{B,S}[F]&=\mathfrak X_{B,T_1}[F]
 =\mathfrak X_{B,T_2}[F]=\mathfrak X_{B,T_3}[F]
 =\mathfrak X_{B,T_4}[F]=0 .
 \label{eq:tensor-xgamma-postborel}
\end{align}
The formerly quoted nonzero $B$ kernels resulted from multiplying this
second propagator term by a free numerator that belongs only to the
$\sigma_{\alpha\beta}$ term.  The corresponding
$\sigma_{\alpha\beta}$ pieces contain no numerator virtuality and are simply
$E_QF_{\rm eff}(a_0)$ times the trace coefficients
\begin{equation}
 \begin{array}{c|cc}
  &A&B\\ \hline
  \mathcal S\ {\rm or}\ \mathcal S_\gamma
    &8&8m_Q/d\\
  \mathcal T_4\ {\rm or}\ \mathcal T_4^\gamma
    &-16\ii&-16\ii m_Q/d
 \end{array}.
 \label{eq:sigma-postborel-coefficients}
\end{equation}
The following operator-level assembly supplies every Fierz, charge, DA and
support weight rather than leaving them implicit.

\paragraph{Explicit operator line projections.}
For an operator-weighted DA $\mathcal F(\alpha_i,v)$, the fraction entering
the heavy denominator is
\begin{equation}
 a(\alpha_i,v)=\alpha_{\bar q}+v\alpha_g .
\end{equation}
Define
\begin{align}
 \mathcal J_z[\mathcal F]={}&
 \int_0^z dv\int_0^{(1-z)/(1-v)}d\alpha_g\,
 \mathcal F(1-z-(1-v)\alpha_g,z-v\alpha_g,\alpha_g,v)
 \nonumber\\
 &+\int_z^1dv\int_0^{z/v}d\alpha_g\,
 \mathcal F(1-z-(1-v)\alpha_g,z-v\alpha_g,\alpha_g,v).
 \label{eq:operator-line-projection}
\end{align}
This is the explicit result of the
$\delta[z-\alpha_{\bar q}-v\alpha_g]$ integration.

The $1/(q\cdot x)$ terms require one further support density.  Performing
Rohrwild's auxiliary-$u$ delta integration before the Borel transformation
gives
\begin{align}
 \mathcal L_z[X]={}&
 \int_0^{1-z}d\alpha_q\int_0^{\alpha_q}d\alpha_q'
 \int_0^z d\alpha_{\bar q}\,
 \frac{z-\alpha_{\bar q}}{(1-\alpha_q-\alpha_{\bar q})^2}
 X(\alpha_q,\alpha_{\bar q},1-\alpha_q-\alpha_{\bar q})
 \nonumber\\
 &-\int_0^{1-z}d\alpha_q\int_z^{1-\alpha_q}d\beta\,
 \frac{z}{\beta^2}\int_0^\beta d\alpha_{\bar q}'\,
 X(\alpha_q,\alpha_{\bar q}',1-\alpha_q-\alpha_{\bar q}')
 \nonumber\\
 &+\int_0^{1-z}d\alpha_q\,\frac{z}{(1-\alpha_q)^2}
 \int_0^{\alpha_q}d\alpha_q'\int_0^{1-\alpha_q'}d\alpha_{\bar q}\,
 X(\alpha_q',\alpha_{\bar q},1-\alpha_q'-\alpha_{\bar q}).
 \label{eq:explicit-P-line-density}
\end{align}
The lower limit $\beta=z$ in the second line is fixed by double-Borel
support.  Extending it to zero, as in the printed symmetric appendix, creates
a spurious logarithmic divergence.

\paragraph{Complete gluonic contribution.}
All DAs in the following operator weights have arguments
$(\alpha_q,\alpha_{\bar q},\alpha_g)$:
\begin{align}
 \mathcal F^g_{A,\sigma}
 &=\mathcal S+\widetilde{\mathcal S}
   -\mathcal T_1-\mathcal T_2+\mathcal T_3+\mathcal T_4,\nonumber\\
 \mathcal F^g_{A,x\gamma}
 &=2v(-\mathcal S-\mathcal T_3+\mathcal T_2),\nonumber\\
 \mathcal F_A^g&=\mathcal F^g_{A,\sigma}
                 +\mathcal F^g_{A,x\gamma},&
 \mathcal F_B^g&=\frac{m_Q}{d}\mathcal F^g_{A,\sigma}.
 \label{eq:operator-weights-gluonic}
\end{align}
Thus
\begin{equation}
 \boxed{\widehat T_A^{\rm 3p,g}
 =e_s\langle\bar ss\rangle E_Q
 \mathcal J_{a_0}[\mathcal F_A^g],
 \qquad
 \widehat T_B^{\rm 3p,g}
 =e_s\langle\bar ss\rangle E_Q
 \mathcal J_{a_0}[\mathcal F_B^g].}
 \label{eq:complete-gluonic-three-particle}
\end{equation}
In particular,
$\mathcal J_{a_0}[\mathcal F_A^g]=\mathcal I_{\mathcal F_1}(u_0)$
for $a_0=1-u_0$.

\paragraph{Complete electromagnetic contribution.}
The direct weights are
\begin{equation}
 \mathcal F_A^\gamma=(1-2v)\mathcal S_\gamma-\mathcal T_4^\gamma,
 \qquad
 \mathcal F_B^\gamma=\frac{m_Q}{d}
 (\mathcal S_\gamma-\mathcal T_4^\gamma).
 \label{eq:operator-weights-em}
\end{equation}
The exact Borel treatment of the $\mathbb P$ numerator replaces the legacy
pole-mass expression according to
\begin{equation}
 \frac{2(m_P^2-M_i^2)}{M^2}\,
 \mathcal L_{a_0}[\mathcal T_4^\gamma]
 \quad\longrightarrow\quad
 2\left.\frac{d}{dz}\mathcal L_z[\mathcal T_4^\gamma]\right|_{z=a_0}.
 \label{eq:P-functional-borel-replacement}
\end{equation}
Writing $\Delta E_Q=E_Q-E_0$, the post-Borel invariants are
\begin{align}
 \boxed{\widehat T_A^{\rm 3p,\gamma}}
 &=e_Q\langle\bar ss\rangle\Delta E_Q
 \left\{\mathcal J_{a_0}[\mathcal F_A^\gamma]
 +2\left.\frac{d}{dz}\mathcal L_z[\mathcal T_4^\gamma]
 \right|_{z=a_0}\right\},
 \nonumber\\
 \boxed{\widehat T_B^{\rm 3p,\gamma}}
 &=e_Q\langle\bar ss\rangle\Delta E_Q\frac{m_Q}{d}
 \left\{\mathcal J_{a_0}[
 \mathcal S_\gamma-\mathcal T_4^\gamma]
 +2\left.\frac{d}{dz}\mathcal L_z[\mathcal T_4^\gamma]
 \right|_{z=a_0}\right\}.
 \label{eq:complete-em-three-particle}
\end{align}
Equations~\eqref{eq:operator-line-projection},
\eqref{eq:explicit-P-line-density}, and
\eqref{eq:complete-em-three-particle} display every integration limit,
operator weight and Borel derivative.  No physical mass and no ordinary local
condensate enters a QCD-side numerator.

For completeness, the axial perturbative density is
\begin{align}
 \rho_i(s;m_i,m_j,e_i)={}&-\frac{3e_i}{8\pi^2}
 \Bigg[
 2m_i\ln\frac{s-m_j^2+m_i^2-\lambda^{1/2}(s,m_j^2,m_i^2)}
 {s-m_j^2+m_i^2+\lambda^{1/2}(s,m_j^2,m_i^2)}
 \nonumber\\
 &+(m_j-m_i)\frac{m_j^2-m_i^2-s}{s^2}
 \lambda^{1/2}(s,m_j^2,m_i^2)\Bigg],\nonumber\\
 \rho_A^P(s)={}&
 \rho_s(s;m_s,m_Q,e_s)-\rho_Q(s;m_Q,m_s,e_Q).
 \label{eq:axial-perturbative-density}
\end{align}

\subsection{Local-scheme separation}

The alternative Colangelo axial comparison adds the ordinary local term
\begin{equation}
 \widehat T_A^{\rm Col,local}
 =e_QE_Q\langle\bar ss\rangle
 \left[1-\frac{m_Qm_s}{M^2}
 +\frac{m_s^2}{2M^2}\left(1+\frac{m_Q^2}{M^2}\right)\right].
 \label{eq:colangelo-local-alternative}
\end{equation}
This term is absent from $\widehat T_A^{\rm RW}$ and
$\widehat T_B^{\rm RW}$.  It must never be added to the electromagnetic
three-particle terms in Eq.~\eqref{eq:complete-em-three-particle}.

\begin{table}[H]
 \centering
 \small
 \begin{tabular}{p{0.34\textwidth}p{0.25\textwidth}p{0.25\textwidth}}
  \toprule
  Contribution & Rohrwild paper default & Colangelo comparison\\
  \midrule
  Ordinary local condensate in transition LCSR & absent & retained for $J_A$\\
  $\mathcal S_\gamma,\mathcal T_4^\gamma$ & nonlocal; explicit post-Borel kernels & not added separately\\
  Gluonic photon DAs & nonlocal & included through $\mathcal F_1$\\
  Local condensates in two-point $f_i$ sum rules & required & required\\
  Local plus electromagnetic three-particle term & forbidden double counting & forbidden double counting\\
  \bottomrule
 \end{tabular}
 \caption{Separation of the local-condensate prescriptions.}
 \label{tab:local-scheme-comparison}
\end{table}

\subsection{Exact pure-axial reduction to Colangelo et al.}

Colangelo, De Fazio, and Ozpineci use the single axial current
$J^A_\mu=\bar Q\gamma_\mu\gamma_5s$ in their correlator.  This is the
Hermitian-conjugate orientation of our $A$ current and has the same invariant
amplitude.  In the convention of Eq.~\eqref{eq:state-rotation}, their
$D_{s1}(2460)$ calculation is obtained from the low channel by choosing
\begin{equation}
  \boxed{\theta=\frac{\pi}{2}=90^\circ},
  \qquad J_1=J_A,\qquad J_2=-J_B.
  \label{eq:colangelo-angle}
\end{equation}
At this point
\begin{equation}
  \widehat T_1=\widehat T_A,\qquad f_1=f_A,\qquad G_1=G_A.
  \label{eq:pure-axial-identities}
\end{equation}
For the equal-Borel convention
$\mathcal M_1^2=\mathcal M_2^2=2M^2$,
Eq.~\eqref{eq:mixed-coupling-from-ope} reduces to
\begin{equation}
  \left.G_1\right|_{\theta=\pi/2}
  =\frac{\exp[(M_1^2+m_P^2)/(2M^2)](m_Q+m_s)}
  {M_1f_A\,m_P^2f_P}\,\widehat T_A^{\rm OPE}.
  \label{eq:colangelo-formula-limit}
\end{equation}
For $Q=c$, $P=D_s$, and $M_1=m_{D_{s1}}$, the right-hand side is exactly the
prefactor and axial OPE appearing in Eq.~(4.4) of
\href{https://arxiv.org/abs/hep-ph/0505195}{Colangelo, De Fazio, and
Ozpineci}.  Therefore the relation is analytic, not merely a numerical
similarity.  Choosing $\theta=0$ instead places the pure axial current in the
second row of Eq.~\eqref{eq:state-rotation}; it swaps the low/high labels and
is not our $D_{s1}(2460)$ convention.

The numerical pure-$A$ comparison is shown in
Table~\ref{tab:colangelo-projection-comparison}.
\begin{table}[H]
  \centering
  \small
  \begin{tabular}{lll}
    \toprule
    Calculation & $G\,[\mathrm{GeV}^{-1}]$ & $\Gamma\,[\mathrm{keV}]$\\
    \midrule
    Colangelo et al., single $J_A$
      & $-0.37$ to $-0.29$ & $19$ to $29$\\
    This project, Stage 2 at $\theta=90^\circ$
      & $-0.3978$ to $-0.3124$ & $20.51$ to $33.26$\\
    Interval overlap
      & $-0.3700$ to $-0.3124$ & $20.51$ to $29.00$\\
    \bottomrule
  \end{tabular}
  \caption{Pure-axial benchmark.  The remaining displacement is caused by the
  input set and convention choices, not by the state projector.}
  \label{tab:colangelo-projection-comparison}
\end{table}
No intermediate angle is fitted here.  Formally, at fixed residues,
$G_1(\theta)=a\sin\theta+b\cos\theta$ with
$a=f_A^{(1)}G_A^{(1)}/f_1$ and
$b=f_B^{(1)}G_B^{(1)}/f_1$.  A physical angle extraction requires the
complete tensor-current transition amplitude.  The angle-dependent physical
residue $f_1(\theta)$ is now supplied independently by the explicit
$AA/AB/BB$ two-point matrix below.

\subsection{Decay constants from the rotated two-point correlator}

Introduce the transverse two-point matrix
\begin{equation}
  \Pi^{\rm basis}=
  \begin{pmatrix}\Pi_{AA}&\Pi_{AB}\\\Pi_{AB}&\Pi_{BB}\end{pmatrix}.
\end{equation}
The physical matrix is
$\Pi^{\rm phys}=R\Pi^{\rm basis}R^T$, so that
\begin{align}
  \Pi_{11}
  &=\sin^2\theta\,\Pi_{AA}+\cos^2\theta\,\Pi_{BB}
    +2\sin\theta\cos\theta\,\Pi_{AB},
  \label{eq:pi11-mixed}\\
  \Pi_{22}
  &=\cos^2\theta\,\Pi_{AA}+\sin^2\theta\,\Pi_{BB}
    -2\sin\theta\cos\theta\,\Pi_{AB},
  \label{eq:pi22-mixed}\\
  \Pi_{12}
  &=\sin\theta\cos\theta(\Pi_{AA}-\Pi_{BB})
    +(\cos^2\theta-\sin^2\theta)\Pi_{AB}.
  \label{eq:pi12-mixed}
\end{align}
This is already the useful paper-level explicit form.  If each basis
correlator is decomposed as
\begin{equation}
  \widehat\Pi_{XY}=
  \int_{s_{\rm th}}^{s_0}ds\,\ee^{-s/\mathcal M^2}
  \rho_{XY}^{\rm pert}(s)
  +\sum_d\widehat\Pi_{XY}^{(d)},
  \qquad XY=AA,BB,AB,
  \label{eq:twopoint-basis-ope}
\end{equation}
All three perturbative densities are explicit.  With
\begin{equation}
 d=m_Q+m_s,\qquad
 \lambda_2(s)=[s-(m_Q+m_s)^2][s-(m_Q-m_s)^2],
\end{equation}
they are
\begin{align}
 \boxed{\rho_{AA}(s)}
 &=\frac{\sqrt{\lambda_2(s)}}{8\pi^2}
 \left[
 2-\frac{m_Q^2+m_s^2+6m_Qm_s}{s}
 -\frac{(m_Q^2-m_s^2)^2}{s^2}
 \right],
 \label{eq:rho-AA-explicit}\\
 \boxed{\rho_{AB}(s)=\rho_{BA}(s)}
 &=\frac{\sqrt{\lambda_2(s)}}{8\pi^2}
 \frac{3(m_s-m_Q)(d^2-s)}{ds},
 \label{eq:rho-AB-explicit}\\
 \boxed{\rho_{BB}(s)}
 &=\frac{\sqrt{\lambda_2(s)}}{8\pi^2}
 \frac{-(d^2-s)
 (2m_s^2-4m_Qm_s+2m_Q^2+s)}
 {d^2s}.
 \label{eq:rho-BB-explicit}
\end{align}
The $AA$ expression is algebraically identical to the Colangelo
axial-current density, and $\rho_{AB}=\rho_{BA}$ is an exact trace check.

At the common LO+$d=3$+$d=5$ truncation, the remaining Borel matrix is also
explicit:
\begin{equation}
 \widehat\Pi(\mathcal M^2,s_0)=
 \int_{d^2}^{s_0}ds\,e^{-s/\mathcal M^2}
 \begin{pmatrix}\rho_{AA}&\rho_{AB}\\
                 \rho_{AB}&\rho_{BB}\end{pmatrix}
 +e^{-m_Q^2/\mathcal M^2}
 \left[C_3(\mathcal Q_0+\mathcal Q_1+\mathcal Q_2)
       +C_5\mathcal Q_5\right],
 \label{eq:complete-twopoint-matrix}
\end{equation}
where $C_3=\langle\bar ss\rangle$,
$C_5=\langle\bar s g_s\sigma\!\cdot\!G s\rangle$, and
\begin{align}
 \mathcal Q_0&=
 \begin{pmatrix}
 m_Q&m_Q^2/d\\m_Q^2/d&m_Q^3/d^2
 \end{pmatrix},
 \nonumber\\
 \mathcal Q_1&=
 \begin{pmatrix}
 -\dfrac{m_sm_Q^2}{2\mathcal M^2}&
 \dfrac{m_sm_Q(1-m_Q^2/\mathcal M^2)}{2d}\\
 \dfrac{m_sm_Q(1-m_Q^2/\mathcal M^2)}{2d}&
 \dfrac{m_sm_Q^2(1-m_Q^2/(2\mathcal M^2))}{d^2}
 \end{pmatrix},
 \nonumber\\
 \mathcal Q_2&=
 \begin{pmatrix}
 \dfrac{m_s^2m_Q^3}{2\mathcal M^4}&
 \dfrac{m_s^2(m_Q^4/\mathcal M^4-m_Q^2/\mathcal M^2-1)}{2d}\\
 \dfrac{m_s^2(m_Q^4/\mathcal M^4-m_Q^2/\mathcal M^2-1)}{2d}&
 \dfrac{m_s^2(m_Q^5/(2\mathcal M^4)-m_Q^3/\mathcal M^2)}{d^2}
 \end{pmatrix},
 \nonumber\\
 \mathcal Q_5&=
 \begin{pmatrix}
 -\dfrac{m_Q^3}{4\mathcal M^4}&
 -\dfrac{m_Q^4/\mathcal M^4-m_Q^2/\mathcal M^2-1}{4d}\\
 -\dfrac{m_Q^4/\mathcal M^4-m_Q^2/\mathcal M^2-1}{4d}&
 -\dfrac{m_Q^5/(4\mathcal M^4)-m_Q^3/(2\mathcal M^2)}{d^2}
 \end{pmatrix}.
 \label{eq:twopoint-local-matrices}
\end{align}
For compactness write $s_\theta=\sin\theta$, $c_\theta=\cos\theta$ and define
the scalar projections of each displayed local matrix by
\begin{align}
 q_1^{(n)}
 &=s_\theta^2(\mathcal Q_n)_{AA}
   +c_\theta^2(\mathcal Q_n)_{BB}
   +2s_\theta c_\theta(\mathcal Q_n)_{AB},\nonumber\\
 q_2^{(n)}
 &=c_\theta^2(\mathcal Q_n)_{AA}
   +s_\theta^2(\mathcal Q_n)_{BB}
   -2s_\theta c_\theta(\mathcal Q_n)_{AB},
 \qquad n=0,1,2,5 .
 \label{eq:projected-local-q}
\end{align}
Because every entry of $\mathcal Q_n$ is given in
Eq.~\eqref{eq:twopoint-local-matrices}, Eq.~\eqref{eq:projected-local-q} is a
fully explicit scalar expression, not a new phenomenological input.

The two poles use separate continuum thresholds
$s_0^{(1)}$ and $s_0^{(2)}$, fitted to $M_1$ and $M_2$, respectively.  The
continuum-subtracted quantities actually evaluated are
\begin{align}
  \widehat\Pi_{11}={}&
  \int_{d^2}^{s_0^{(1)}}ds\,\ee^{-s/\mathcal M^2}
  \left[
    s_\theta^2\rho_{AA}(s)
    +c_\theta^2\rho_{BB}(s)
    +2s_\theta c_\theta\rho_{AB}(s)
  \right]
  \nonumber\\
  &+\ee^{-m_Q^2/\mathcal M^2}
  \left\{C_3[q_1^{(0)}+q_1^{(1)}+q_1^{(2)}]
  +C_5q_1^{(5)}\right\},
  \label{eq:pi11-explicit-ope}\\
  \widehat\Pi_{22}={}&
  \int_{d^2}^{s_0^{(2)}}ds\,\ee^{-s/\mathcal M^2}
  \left[
    c_\theta^2\rho_{AA}(s)
    +s_\theta^2\rho_{BB}(s)
    -2s_\theta c_\theta\rho_{AB}(s)
  \right]
  \nonumber\\
  &+\ee^{-m_Q^2/\mathcal M^2}
  \left\{C_3[q_2^{(0)}+q_2^{(1)}+q_2^{(2)}]
  +C_5q_2^{(5)}\right\}.
  \label{eq:pi22-explicit-ope}
\end{align}
The final square-root formulas used in the numerical code are therefore
\begin{align}
  \boxed{f_1(\mathcal M^2,s_0^{(1)},\theta)}
  &=\frac{\ee^{M_1^2/(2\mathcal M^2)}}{M_1}
  \sqrt{\widehat\Pi_{11}(\mathcal M^2,s_0^{(1)},\theta)},
  \label{eq:f1-mixed-sumrule}\\
  \boxed{f_2(\mathcal M^2,s_0^{(2)},\theta)}
  &=\frac{\ee^{M_2^2/(2\mathcal M^2)}}{M_2}
  \sqrt{\widehat\Pi_{22}(\mathcal M^2,s_0^{(2)},\theta)}.
  \label{eq:f2-mixed-sumrule}
\end{align}
The condition $\widehat\Pi_{12}\simeq0$ can be used to diagnose or determine
the preferred mixing angle.  At $\theta=\pi/2$,
Eq.~\eqref{eq:f1-mixed-sumrule} becomes the pure axial-current result
\begin{equation}
  f_1=\frac{\ee^{M_1^2/(2\mathcal M^2)}}{M_1}
  \sqrt{\widehat\Pi_{AA}},
\end{equation}
which is the decay-constant normalization used by Colangelo et al.

\subsection{Current two-point result for $f_1$ and $f_2$}

Equations~\eqref{eq:rho-AA-explicit}--\eqref{eq:f2-mixed-sumrule}
replace the former external-$f_1$/overlap normalization.  At the fixed angle
from our previous study,
\begin{equation}
 \theta_{D_s}=26.6\pm0.6^\circ,
\end{equation}
the fixed-angle, two-threshold Monte Carlo gives
\begin{align}
 f_1&=0.4046[0.3799,0.4291]~\mathrm{GeV},&
 f_2&=0.1677[0.1594,0.1756]~\mathrm{GeV}.
 \label{eq:current-f1-f2-result}
\end{align}
For example, the explicit central Mathematica substitution
\begin{equation}
 \mathcal M^2=2.35~{\rm GeV}^2,\qquad
 s_0^{(1)}=7.925~{\rm GeV}^2,\qquad
 s_0^{(2)}=9.975~{\rm GeV}^2,\qquad
 \theta=26.6^\circ
\end{equation}
with $m_c=1.27~{\rm GeV}$, $m_s=0.093~{\rm GeV}$,
$\langle\bar ss\rangle=0.8[-(0.24~{\rm GeV})^3]$ and
$m_0^2=0.8~{\rm GeV}^2$ gives
\begin{equation}
 f_1=0.4005893~{\rm GeV},\qquad
 f_2=0.1666704~{\rm GeV}.
 \label{eq:central-f1-f2-explicit}
\end{equation}
The quoted values in Eq.~\eqref{eq:current-f1-f2-result} are the medians and
16--84 percentile intervals after varying the QCD inputs, Borel parameter,
threshold fits and the external angle.
No decay constant is externally anchored in this result.  Diagonalizing the
same truncated matrix gives
$\theta_{\rm matrix}=30.63[29.94,31.29]^\circ$; this is retained as an
OPE-truncation diagnostic rather than replacing the dedicated external angle.
The Python and Mathematica fixed-point residues agree below
$3\times10^{-8}\,\mathrm{GeV}$.  Local condensates are required in this
ordinary two-point sum rule; their presence here does not reintroduce a local
condensate into the external-photon transition LCSR.

Finally, once $G_i$ is normalized with the corresponding $f_i$, the decay
width is
\begin{equation}
  \Gamma(H_i\to P\gamma)
  =\frac{\alpha_{\rm em}}{3}|G_i|^2k_{\gamma,i}^3,
  \qquad
  k_{\gamma,i}=\frac{M_i^2-m_P^2}{2M_i}.
  \label{eq:mixed-radiative-width}
\end{equation}

\section{Inverse-denominator reduction}

The projected expressions contain $k^2$, $k\cdot p$, and $k\cdot q$.  For the
heavy-emission topology these are replaced by
\begin{align}
  k^2&=d_{H2}+m_Q^2,\\
  k\cdot q&=\frac{d_{H1}-d_{H2}}{2},\\
  k\cdot p&=\frac{d_{H2}+m_Q^2+p^2-m_s^2-d_{H3}}{2}.
\end{align}
For strange emission,
\begin{align}
  k^2&=d_{S1}+m_Q^2,\\
  k\cdot p&=\frac{d_{S1}+m_Q^2+p^2-m_s^2-d_{S2}}{2},\\
  k\cdot q&=p\cdot q+\frac{d_{S2}-d_{S3}}{2}.
\end{align}

Setting all three inverse denominators of a topology to zero extracts the
coefficient of the genuine three-denominator triangle.  Terms containing one
or more $d_i$ cancel propagators and are retained separately in the full
machine output.

The axial triangle cores are
\begin{align}
  N^{A,\mathrm{core}}_H
  &=4\ii m_Q+
    \frac{4\ii m_Q^3-4\ii m_Q^2m_s}{p\cdot q},\\
  N^{A,\mathrm{core}}_s
  &=-4\ii m_s+
    \frac{4\ii m_Qm_s^2-4\ii m_s^3}{p\cdot q}.
\end{align}
The tensor-current cores are
\begin{align}
  N^{B,\mathrm{core}}_H
  &=\frac{8\ii m_Q^2-4\ii m_Qm_s}{m_Q+m_s}
    +\frac{4\ii m_Q^2p^2}{(m_Q+m_s)(p\cdot q)},\\
  N^{B,\mathrm{core}}_s
  &=\frac{-4\ii m_Qm_s+8\ii m_s^2}{m_Q+m_s}
    +\frac{4\ii m_s^2p^2}{(m_Q+m_s)(p\cdot q)}.
\end{align}
These expressions refer only to the numerator traces of
Eqs.~\eqref{eq:heavy-trace}--\eqref{eq:strange-trace}; the charge, color,
propagator, QED-vertex, and global Wick factors must still be attached.

\section{Feynman parameterization of the hard triangles}

For three denominators,
\begin{equation}
  \frac{1}{D_1D_2D_3}
  =2\int_0^1dx\int_0^{1-x}dy\,
  \frac{1}{[xD_1+yD_2+(1-x-y)D_3]^3}.
  \label{eq:feynman-identity}
\end{equation}
Equivalently one may use $x+y+z=1$.

\subsection{Heavy emission}

Assign $x,y,z$ to $d_{H1},d_{H2},d_{H3}$.  Direct completion of the square
gives
\begin{equation}
  xd_{H1}+yd_{H2}+zd_{H3}
  =(k+xq-zp)^2-\Delta_H,
\end{equation}
where
\begin{equation}
  \boxed{\Delta_H=(x+y)m_Q^2+zm_s^2-xz p^{\prime2}-yzp^2.}
  \label{eq:deltaH}
\end{equation}

\subsection{Strange emission}

Assign $x,y,z$ to $d_{S1},d_{S2},d_{S3}$.  Then
\begin{equation}
  xd_{S1}+yd_{S2}+zd_{S3}
  =[k-yp-z(p+q)]^2-\Delta_s,
\end{equation}
with
\begin{equation}
  \boxed{\Delta_s=xm_Q^2+(y+z)m_s^2-xyp^2-xz p^{\prime2}.}
  \label{eq:deltaS}
\end{equation}

The independent pairwise-invariant formula
\begin{equation}
  \Delta=\sum_i x_i m_i^2
  -\sum_{i<j}x_ix_j(r_i-r_j)^2
\end{equation}
reproduces Eqs.~\eqref{eq:deltaH} and \eqref{eq:deltaS}.  The notebook checks
both the completed-square and pairwise forms:
\begin{lstlisting}
deltaH = (x + y) mQ^2 + z ms^2 - x z pPrime2 - y z p2;
shiftH = x q - z p;
heavyResidual = simplexReduce[
  x dH1 + y dH2 + z dH3
  - (SP[k + shiftH, k + shiftH] - deltaH)
];

deltaS = x mQ^2 + (y + z) ms^2 - x y p2 - x z pPrime2;
shiftS = -y p - z (p + q);
strangeResidual = simplexReduce[
  x dS1 + y dS2 + z dS3
  - (SP[k + shiftS, k + shiftS] - deltaS)
];
\end{lstlisting}
Both residuals vanish exactly.

\section{Soft two-particle Fierz traces}

After contracting the heavy quark and Fierz rearranging the strange bilocal,
the Dirac factor is
\begin{equation}
  T_i^X=\Tr_D[\Gamma_P(\slashedk+m_Q)\Gamma^X_\mu\Gamma_i].
\end{equation}
The notebook evaluates the standard scalar, pseudoscalar, vector, axial, and
tensor basis.  The nonzero results for the axial current are
\begin{align}
  T^A_{1}&=4\ii k_\mu,\\
  T^A_{\gamma_\rho}&=-4\ii m_Q g_{\mu\rho},\\
  T^A_{\sigma_{\rho\lambda}}
  &=4g_{\lambda\mu}k_\rho-4g_{\mu\rho}k_\lambda.
\end{align}
For the tensor current,
\begin{align}
  T^B_{\gamma_\rho}
  &=\frac{4\ii}{m_Q+m_s}
    [k_\mu p^{\prime}_\rho-g_{\mu\rho}(k\cdot p^{\prime})],\\
  T^B_{\gamma_\rho\gamma_5}
  &=-\frac{4}{m_Q+m_s}\epsilon_{\mu\rho k p^{\prime}},\\
  T^B_{\sigma_{\rho\lambda}}
  &=\frac{4m_Q}{m_Q+m_s}
    [g_{\lambda\mu}p^{\prime}_\rho-g_{\mu\rho}p^{\prime}_\lambda].
\end{align}
All omitted entries vanish.  Cyclically rotating the matrix chain gives the
same tensor- and vector-basis traces for both currents.

\clearpage
\section{Implemented OPE terms: exact status}

The Mathematica restart now contains the complete analytic operator-weighted
post-Borel assembly.  The Python and Mathematica numerical layers now use the
same corrected off-shell expressions; the older pole-mass numerator routines
remain only as explicitly labeled historical diagnostics.  The present status
is:

\begin{table}[H]
  \centering
  \small
  \begin{tabular}{p{0.27\textwidth}p{0.27\textwidth}p{0.32\textwidth}}
    \toprule
    OPE contribution & Axial current $A$ & Tensor current $B$\\
    \midrule
    Perturbative photon from $Q$ and $s$
      & Published spectral density implemented numerically; corrected traces
        independently reproduced
      & Corrected traces and triangle cores derived; diagonal
        double-dispersion spectral mapping implemented independently in
        Python and Mathematica\\
    Local $e_Q\langle\bar ss\rangle$ term
      & Excluded in the Rohrwild result; retained only in the Colangelo
        comparison with $m_s,m_s^2$ corrections
      & Excluded in the Rohrwild result; no Colangelo $J_B$ term invented\\
    Electromagnetic three-particle DAs,
      $\mathcal S_\gamma,\mathcal T_4^\gamma$
      & Complete support-corrected operator convolution is explicit
      & Complete support-corrected operator convolution is explicit; the
        former pole-mass numerator substitution is withdrawn\\
    Twist-2 two-particle DA,
      $\chi\langle\bar ss\rangle\phi_\gamma$
      & Included in the exact numerical invariant
      & Exact factor $m_Q/(m_Q+m_s)$ included\\
    Twist-4 two-particle DAs,
      $\mathcal A$ and $\mathcal B$
      & Explicit in the Rohrwild combination $\mathcal A+2\mathcal B$
      & Exact factor $m_Q/(m_Q+m_s)$ relative to the axial piece\\
    Twist-3 two-particle DA,
      $f_{3\gamma}\overline\psi^{(V)}$
      & Explicit; vanishes at the symmetric Borel point
      & Exact post-Borel derivative expression in
        Eq.~\eqref{eq:tensor-vector-explicit-borel}\\
    Three-particle photon DAs,
      $\mathcal S,\widetilde{\mathcal S},\mathcal T_{1\ldots4}$
      & Full operator-weighted $\mathcal F_1$ line projection is explicit
      & Full operator-weighted line projection is explicit; the
        $x_\alpha\gamma_\beta$ trace vanishes\\
    Mixed condensate
      $\langle\bar s g_s\sigma\cdot Gs\rangle$
      & Not included in the radiative three-point sum rule
      & Not included\\
    Vacuum gluon condensate $\langle G^2\rangle$
      & Not included in the radiative three-point sum rule
      & Not included\\
    \bottomrule
  \end{tabular}
  \caption{Term-by-term status of the completed analytic and numerical
  transition calculation.}
  \label{tab:ope-status}
\end{table}

\subsection{Completed numerical checkpoint}

At $a_0=1/2$ the independently evaluated DA convolutions are
\begin{align}
\mathcal J[\mathcal F^g_{A,\sigma}]&=-0.1562500000,&
\mathcal J[\mathcal F^g_{A,x\gamma}]&=+0.2265798500,\nonumber\\
\mathcal J[\mathcal F^g_A]&=+0.0703298500,&
\mathcal J[\mathcal F^\gamma_A]&=0,\nonumber\\
\mathcal J[\mathcal S_\gamma-\mathcal T_4^\gamma]&=1.5625000000,&
\mathcal L[\mathcal T_4^\gamma]&=0.1111834583,\nonumber\\
\mathcal L'[\mathcal T_4^\gamma]&=0.4098669166.&&
\end{align}
For the charm reference point
$M^2=3.75~\GeV^2$, $s_0=8.0~\GeV^2$ and the lattice photon
normalization, the totals are
\begin{equation}
\widehat T_A=-0.0311671290~\GeV^3,\qquad
\widehat T_B=-0.0701320395~\GeV^3.
\end{equation}
With $\theta=26.6^\circ$ and the central projected residues this gives
\begin{align}
G_1&=-0.4113423853~\GeV^{-1},&
\Gamma_1&=35.5665446~\keV,\nonumber\\
G_2&=+0.0464986762~\GeV^{-1},&
\Gamma_2&=0.6709364~\keV.
\end{align}
These are single-point regression values rather than final-window medians.
For the direct bottom two-point reference
$M_{\rm 2pt}^2=7.00~\GeV^2$,
$s_0^{(1)}=44.35~\GeV^2$, and
$s_0^{(2)}=43.025~\GeV^2$, the independent Mathematica evaluation gives
\begin{equation}
 f_1^{(B_s)}=0.5381107~\GeV,\qquad
 f_2^{(B_s)}=0.0888787~\GeV,
\end{equation}
in agreement with Python.  The direct $B_s$ Monte Carlo gives
$f_1^{(B_s)}=0.5358[0.5096,0.5586]~\GeV$ and
$f_2^{(B_s)}=0.08896[0.08496,0.09278]~\GeV$.
The cross-code transition comparison contains 145 common numerical entries and
passes the $10^{-8}$ absolute-or-relative tolerance.  The full fixed-seed
uncertainty scan gives
\begin{align}
\Gamma[D_{s1}(2460)\to D_s\gamma]
 &=34.96[28.06,43.53]~\keV,\nonumber\\
\Gamma[D_{s1}(2536)\to D_s\gamma]
 &=0.988[0.092,3.555]~\keV,\nonumber\\
\Gamma[B_{s1}(5750)\to B_s\gamma]
 &=11.34[9.144,14.06]~\keV,\nonumber\\
\Gamma[B_{s1}(5830)\to B_s\gamma]
 &=8.702[6.474,11.23]~\keV.
\end{align}

\section{Verification summary and present boundary}

\begin{table}[H]
  \centering
  \renewcommand{\arraystretch}{1.12}
  \begin{tabular}{p{0.52\textwidth}p{0.09\textwidth}p{0.25\textwidth}}
    \toprule
    Check & Result & Exact tests\\
    \midrule
    Currents, Wick signs, traces, E1 cores, soft trace cyclicity
      & PASS & 19/19\\
    Standard vertex routing
      & PASS & 4/4 momenta\\
    Heavy- and strange-line Ward identities
      & PASS & 2/2\\
    Shifted and pairwise Feynman denominators
      & PASS & 4/4\\
    State rotation, OPE assembly, Colangelo limit, decay constant, and width
      & PASS & 15/15\\
    Notebook construction and corrected numerical section
      & PASS & 53 cells; block run\\
    Python--Mathematica transition regression
      & PASS & 145/145 quantities\\
    Direct $B_s$ two-point central regression
      & PASS & $f_1,f_2$, thresholds, pole/OPE diagnostics\\
    \bottomrule
  \end{tabular}
  \caption{Status of the independent Mathematica/FeynCalc restart.}
\end{table}

The $B_s$ channel uses the same symbolic expressions with $m_Q\to m_b$ and
the appropriate heavy-quark charge.  Its direct $AA/AB/BB$ two-point
normalization is implemented independently in
\path{scripts/mathematica_twopoint_bs1_matrix_check.wl}; no
Pullin--Zwicky decay constant or overlap closure enters the final result.

The completed symbolic and numerical stages now include the exact
two-particle vector-DA Borel reduction, the corrected
$x_\alpha\gamma_\beta$ kernels, the full operator-weighted gluonic and
electromagnetic line projections, and the physical-state uncertainty scan.
Useful higher-order follow-up tasks are:
\begin{enumerate}
  \item independently rederive the shifted hard-loop double discontinuity
    directly from Eqs.~\eqref{eq:deltaH}--\eqref{eq:deltaS}, as a check beyond
    the present published spectral-density implementation;
  \item add perturbative $\alpha_s$ corrections consistently to the transition
    and two-point sectors;
  \item extend the common OPE truncation to the omitted higher-dimensional
    terms.
\end{enumerate}
Python results that use the former physical-residue numerator shortcut remain
historical diagnostics and are not used in the quoted tables.

\appendix
\section{Reproduction commands}

From \texttt{papers/Ds\_Bs\_gamma/}, run
\begin{lstlisting}[language={},basicstyle=\ttfamily\scriptsize]
/Applications/Wolfram.app/Contents/MacOS/WolframKernel -noprompt \
  -script scripts/mathematica_current_wick_trace_audit.wl

/Applications/Wolfram.app/Contents/MacOS/WolframKernel -noprompt \
  -script scripts/mathematica_correlator_routing_audit.wl

/Applications/Wolfram.app/Contents/MacOS/WolframKernel -noprompt \
  -script scripts/mathematica_hard_loop_parameterization.wl

/Applications/Wolfram.app/Contents/MacOS/WolframKernel -noprompt \
  -script scripts/mathematica_state_mixing_projection.wl

/Applications/Wolfram.app/Contents/MacOS/WolframKernel -noprompt \
  -script scripts/step13_soft_vector_da_borel.wl

/Applications/Wolfram.app/Contents/MacOS/WolframKernel -noprompt \
  -script scripts/step14_explicit_double_borel_forms.wl

/Applications/Wolfram.app/Contents/MacOS/WolframKernel -noprompt \
  -script scripts/step15_complete_three_particle_borel.wl

python3 scripts/corrected_transition_central_python.py

/Applications/Wolfram.app/Contents/MacOS/WolframKernel -noinit -noprompt \
  -script scripts/mathematica_corrected_transition_numerics.wl

python3 scripts/compare_corrected_transition_numerics.py

python3 scripts/final_window_monte_carlo.py
\end{lstlisting}
To regenerate the notebook from its source cells, run
\begin{lstlisting}[language={},basicstyle=\ttfamily\scriptsize]
/Applications/Wolfram.app/Contents/MacOS/WolframKernel -noprompt \
  -script scripts/build_dsbs_symbolic_notebook.wl

/Applications/Wolfram.app/Contents/MacOS/WolframKernel -noinit -noprompt \
  -script scripts/check_dsbs_symbolic_notebook.wl
\end{lstlisting}

\end{document}
